% 한글 초록
\begin{sloppypar}
바이러스 게놈의 변이 핫스팟 탐지는 변이주 감시와 백신 설계에 필수적이나, 탐지 방법론을 체계적으로 비교·평가할 수 있는 벤치마크 프레임워크가 부재하였다.
기존 연구들은 단일 병원체·단일 탐지 방법에 의존하여, 어떤 점수화(scoring)와 탐지(detection) 조합이 어떤 병원체에 적합한지 알 수 없었다.

본 연구에서는 이 문제를 해결하기 위해 MutBench를 제안한다.
MutBench는 변이 핫스팟 탐지를 위한 최초의 체계적 벤치마크 프레임워크로, 세 가지 핵심 구성요소를 포함한다:
(1)~같은 바이러스 종 내에서 서로 다른 계통에서 독립적으로 반복 출현한 수렴 진화 위치를 양성 선택 기준(Layer~A),
정화 선택 하의 보존 구조 영역을 음성 기준(Layer~B),
Deep Mutational Scanning(DMS) 실험적 적합도 데이터를
기능적 검증(Layer~C)으로 통합하는 3층 ground truth 체계,
(2) $\text{hotspot-score} = (\text{recall} + \text{precision} + \text{stability}) / 3$으로 정의되는 복합 평가 지표,
(3) Stage~1에서 SARS-CoV-2에 대한 심층 분석, Stage~2에서 11개 병원체에 대한 8{,}580회 평가(20개 점수화 공식 $\times$ 39개 탐지 알고리즘 $\times$ 11개 병원체)를 수행하는 2단계 실험 설계이다.

Stage~1에서는 MutClust-Hybrid가 최고 hotspot-score(0.778)를 달성하여,
도메인 특화 사전지식과 데이터 기반 클러스터링의 결합이 효과적임을 확인하였다.
Stage~2의 11개 병원체 대규모 벤치마크에서는 점수화(scoring) 방법과 병원체 간 상호작용이 전체 성능 분산의 29.6\%를 설명하여 최대 분산 원천으로 확인되었다($\omega^2_{\text{scoring} \times \text{pathogen}} = 0.296$). 이는 탐지 알고리즘 패밀리 주효과($\omega^2 = 0.013$)의 22배로, ``어떤 알고리즘을 쓰느냐''보다 ``어떤 생물학적 정보를 쓰느냐''가 성능을 결정하는 핵심 요인임을 입증한다.
Friedman 검정($\chi^2 = 7.69$, $p = 0.990$)에서 상위 20개 조합 간 유의한 차이가 없음을 확인하였으며, LOPO 교차검증에서 11개 병원체 모두 서로 다른 최적 조합을 보였다(일반화 성공률 0/11).

이러한 결과는 최적 탐지 전략이 병원체에 따라 근본적으로 상이하며, 단일 고정 방법이 아닌 병원체별 적응적 접근이 필요함을 통계적으로 입증한다.
6개 정보 카테고리(빈도 기반, MSA 유래, 계통 분석, 구조 기반, AI 기반, 복합)에 걸친 20개 점수화 유형 분석에서 9가지 서로 다른 최적 scoring이 관찰되었다: FUBAR 계통선택 점수(H3N2), 호모플래시(Norovirus), 단백질 언어모델(SARS-CoV-2, RSV), 의미적 변화(EV-A71), 안정성 예측(Rabies), 빈도 기반(HCV, Influenza B, MERS).
또한 탐지된 핫스팟이 실험적으로 확인된 백신 escape 위치와 4.90배 enrichment($p < 0.0001$)를 보여, 공중보건 감시에 대한 실용적 가치를 입증하였다.
다중 정보 통합이 실험적 탐색 범위를 축소함을 확인하였다(H3N2 백신 escape enrichment 4.012배, $p = 0.0023$).

\vspace{1em}
\end{sloppypar}
\noindent\textbf{주요어}: 변이 핫스팟 탐지, 벤치마크 프레임워크, 바이러스 게놈, 복합 평가 지표, 다중 정보 통합 분석, SARS-CoV-2, Deep Mutational Scanning
